\chapter{Приложения}
\section{Приложение №1}

\begin{minted}{Julia}
    begin
        using Quaternions
        using DualNumbers
        using LinearAlgebra
        
        # Вспомогательные функции для кватернионов
        function dot(q1::Quaternion, q2::Quaternion)
            return q1.s * q2.s + q1.v1 * q2.v1 + q1.v2 * q2.v2 + q1.v3 * q2.v3
        end
        
        function cross(q1::Quaternion, q2::Quaternion)
            # Векторное произведение для чисто векторных кватернионов
            return Quaternion(0.0, 
                q1.v2 * q2.v3 - q1.v3 * q2.v2,
                q1.v3 * q2.v1 - q1.v1 * q2.v3,
                q1.v1 * q2.v2 - q1.v2 * q2.v1
            )
        end
        
        function conj(q::Quaternion)
            return Quaternion(q.s, -q.v1, -q.v2, -q.v3)
        end
        
        function norm(q::Quaternion)
            return sqrt(dot(q, q))
        end
        
        struct DualQuaternion
            q::Quaternion{Float64}   # вещественная часть
            qo::Quaternion{Float64}  # дуальная часть
            
            # Конструктор из 4 дуальных чисел
            function DualQuaternion(Q0::Dual, Q1::Dual, Q2::Dual, Q3::Dual)
                q = Quaternion(Q0.value, Q1.value, Q2.value, Q3.value)
                qo = Quaternion(Q0.epsilon, Q1.epsilon, Q2.epsilon, Q3.epsilon)
                new(q, qo)
            end
            
            # Конструктор из двух кватернионов
            DualQuaternion(q::Quaternion, qo::Quaternion) = new(q, qo)
        end
        
        # Арифметические операции
        Base.:+(Q::DualQuaternion, P::DualQuaternion) = DualQuaternion(Q.q + P.q, Q.qo + P.qo)
        Base.:-(Q::DualQuaternion) = DualQuaternion(-Q.q, -Q.qo)
        Base.:-(Q::DualQuaternion, P::DualQuaternion) = DualQuaternion(Q.q - P.q, Q.qo - P.qo)
        
        # Сложение с дуальным числом (только скалярная часть)
        Base.:+(Q::DualQuaternion, d::Dual) = DualQuaternion(
            Quaternion(Q.q.s + d.value, Q.q.v1, Q.q.v2, Q.q.v3),
            Quaternion(Q.qo.s + d.epsilon, Q.qo.v1, Q.qo.v2, Q.qo.v3)
        )
        Base.:+(d::Dual, Q::DualQuaternion) = Q + d
        Base.:-(Q::DualQuaternion, d::Dual) = Q + (-d)
        Base.:-(d::Dual, Q::DualQuaternion) = DualQuaternion(
            Quaternion(d.value - Q.q.s, -Q.q.v1, -Q.q.v2, -Q.q.v3),
            Quaternion(d.epsilon - Q.qo.s, -Q.qo.v1, -Q.qo.v2, -Q.qo.v3)
        )
        
        # Умножение на число
        Base.:*(Q::DualQuaternion, a::Real) = DualQuaternion(a * Q.q, a * Q.qo)
        Base.:*(a::Real, Q::DualQuaternion) = Q * a
        
        # Умножение на дуальное число
        Base.:*(Q::DualQuaternion, d::Dual) = DualQuaternion(
            d.value * Q.q,
            d.value * Q.qo + d.epsilon * Q.q
        )
        Base.:*(d::Dual, Q::DualQuaternion) = Q * d
        
        # Умножение двух дуальных кватернионов
        Base.:*(Q::DualQuaternion, P::DualQuaternion) = DualQuaternion(
            Q.q * P.q,
            Q.q * P.qo + Q.qo * P.q
        )
        
        # Деление
        Base.:/(Q::DualQuaternion, a::Real) = DualQuaternion(Q.q / a, Q.qo / a)
        Base.:/(Q::DualQuaternion, d::Dual) = Q * inv(d)
        
        # Сравнение
        Base.:(==)(Q1::DualQuaternion, Q2::DualQuaternion) = (Q1.q == Q2.q) && (Q1.qo == Q2.qo)
        
        function Base.isapprox(Q1::DualQuaternion, Q2::DualQuaternion; 
                              atol::Real=1e-12, rtol::Real=0)
            return isapprox(Q1.q.s, Q2.q.s; atol, rtol) &&
                   isapprox(Q1.q.v1, Q2.q.v1; atol, rtol) &&
                   isapprox(Q1.q.v2, Q2.q.v2; atol, rtol) &&
                   isapprox(Q1.q.v3, Q2.q.v3; atol, rtol) &&
                   isapprox(Q1.qo.s, Q2.qo.s; atol, rtol) &&
                   isapprox(Q1.qo.v1, Q2.qo.v1; atol, rtol) &&
                   isapprox(Q1.qo.v2, Q2.qo.v2; atol, rtol) &&
                   isapprox(Q1.qo.v3, Q2.qo.v3; atol, rtol)
        end
        
        # Сопряжения
        dualconj(Q::DualQuaternion) = DualQuaternion(Q.q, -Q.qo)
        complexconj(Q::DualQuaternion) = DualQuaternion(conj(Q.q), conj(Q.qo))
        conj(Q::DualQuaternion) = DualQuaternion(conj(Q.q), -conj(Q.qo))
        
        # Нормы и инварианты
        abs2(Q::DualQuaternion) = Dual(dot(Q.q, Q.q), 2 * dot(Q.q, Q.qo))
        invariant(Q::DualQuaternion) = dot(Q.q, Q.qo)
        parameter(Q::DualQuaternion) = invariant(Q) / dot(Q.q, Q.q)
        abs(Q::DualQuaternion) = Dual(norm(Q.q), invariant(Q) / norm(Q.q))
        
        # Конструкторы специальных объектов
        function screw(v::AbstractVector, m::AbstractVector)
            @assert length(v) == 3 && length(m) == 3 "Vectors must have length 3"
            return DualQuaternion(
                Quaternion(0.0, v[1], v[2], v[3]),
                Quaternion(0.0, m[1], m[2], m[3])
            )
        end
        
        function mass_point(w::Real, p::AbstractVector)
            @assert length(p) == 3 "Point must have length 3"
            return DualQuaternion(
                Quaternion(w, 0.0, 0.0, 0.0),
                Quaternion(0.0, p[1], p[2], p[3])
            )
        end
        
        function dq_plane(n::AbstractVector, d::Real)
            @assert length(n) == 3 "Normal vector must have length 3"
            return DualQuaternion(
                Quaternion(0.0, n[1], n[2], n[3]),
                Quaternion(d, 0.0, 0.0, 0.0)
            )
        end
        
        # Скалярное и векторное произведение
        dot(P::DualQuaternion, Q::DualQuaternion) = Dual(
            dot(P.q, Q.q),
            dot(P.qo, Q.q) + dot(P.q, Q.qo)
        )
        
        function cross(P::DualQuaternion, Q::DualQuaternion)
            @assert isapprox(P.q.s, 0.0) && isapprox(P.qo.s, 0.0) && 
                    isapprox(Q.q.s, 0.0) && isapprox(Q.qo.s, 0.0) "Screw product requires pure dual quaternions"
            return DualQuaternion(
                cross(P.q, Q.q),
                cross(P.q, Q.qo) + cross(P.qo, Q.q)
            )
        end
        
        # Формулы для винтового движения
        function RodriguesHamilton(A::DualQuaternion, Θ::Dual)
            Λ0 = cos(Θ/2)
            Λ = sin(Θ/2) * A
            return L -> L + 2Λ0 * cross(Λ, L) + 2 * cross(Λ, cross(Λ, L))
        end
        
        function LineSandwichFormula(A::DualQuaternion, Θ::Dual)
            Λ = cos(Θ/2) + sin(Θ/2) * A
            return L -> Λ * L * complexconj(Λ)
        end
        
        function PointSandwichFormula(A::DualQuaternion, Θ::Dual)
            Λ = cos(Θ/2) + sin(Θ/2) * A
            return P -> Λ * P * conj(Λ)
        end
        
        PlaneSandwichFormula = PointSandwichFormula
        
        # Вспомогательная функция для ортогональных векторов
        function orthogonal_triple(n::AbstractVector)
            @assert length(n) == 3 "Vector must have length 3"
            if n[1] <= n[2] && n[1] <= n[3]
                w = n + [1.0, 0.0, 0.0]
            elseif n[2] <= n[1] && n[2] <= n[3]
                w = n + [0.0, 1.0, 0.0]
            else
                w = n + [0.0, 0.0, 1.0]
            end
            u = normalize(cross(w, n))
            v = normalize(cross(n, u))
            return (v, u)
        end
        
        
    end

    begin
    points = [
        [0.0, 0.0, 0.0],      # A
        [1.0, 0.0, 0.0],      # B
        [1.0, 1.0, 0.0],      # C
        [0.0, 1.0, 0.0],      # D
        [0.5, 0.5, √2 / 2]    # S (вершина)
    ]
    
    function point_to_dualquaternion(p::Vector)
        return DualQuaternion(
            Quaternion(1.0, 0.0, 0.0, 0.0),
            Quaternion(0.0, p[1], p[2], p[3])
        )
    end
    
    function dualquaternion_to_point(dq::DualQuaternion)
        return [dq.qo.v1, dq.qo.v2, dq.qo.v3]
    end
    
    # Грани пирамиды
    base_faces = [[1,2,3,4]]
    side_faces = [
        [1,2,5], [2,3,5], [3,4,5], [4,1,5]
    ]
    end

    function plot_pyramids(original, transformed)
        p = plot(size=(1000, 800), legend=:topleft, 
                xlabel="X", ylabel="Y", zlabel="Z",
                camera=(30, 30),
                xlims=(-2, 2), ylims=(-2, 2), zlims=(-0.2, 2.2))
        
        orig_x = [pt[1] for pt in original]
        orig_y = [pt[2] for pt in original] 
        orig_z = [pt[3] for pt in original]
        
        trans_x = [pt[1] for pt in transformed]
        trans_y = [pt[2] for pt in transformed]
        trans_z = [pt[3] for pt in transformed]

        # Исходная пирамида
        base_x = [orig_x[1], orig_x[2], orig_x[3], orig_x[4], orig_x[1]]
        base_y = [orig_y[1], orig_y[2], orig_y[3], orig_y[4], orig_y[1]]
        base_z = [orig_z[1], orig_z[2], orig_z[3], orig_z[4], orig_z[1]]
        plot!(p, base_x, base_y, base_z, color=:red, linewidth=2, label="Исходная пирамида")
        
        for face in side_faces
            fx = [orig_x[i] for i in face]
            fy = [orig_y[i] for i in face] 
            fz = [orig_z[i] for i in face]
            plot!(p, fx, fy, fz, color=:red, linewidth=2, alpha=0.7, label="")
        end

        # Преобразованная пирамида
        trans_base_x = [trans_x[1], trans_x[2], trans_x[3], trans_x[4], trans_x[1]]
        trans_base_y = [trans_y[1], trans_y[2], trans_y[3], trans_y[4], trans_y[1]]
        trans_base_z = [trans_z[1], trans_z[2], trans_z[3], trans_z[4], trans_z[1]]
        plot!(p, trans_base_x, trans_base_y, trans_base_z, color=:blue, linewidth=2, label="После винтового движения")
        
        for face in side_faces
            fx = [trans_x[i] for i in face]
            fy = [trans_y[i] for i in face]
            fz = [trans_z[i] for i in face]
            plot!(p, fx, fy, fz, color=:blue, linewidth=2, alpha=0.7, label="")
        end

        # Точки
        scatter!(p, orig_x, orig_y, orig_z, color=:red, markersize=6, marker=:circle, label="")
        scatter!(p, trans_x, trans_y, trans_z, color=:blue, markersize=6, marker=:circle, label="")

        # Подписи
        labels = ["A", "B", "C", "D", "S"]
        offset = 0.05
        for i in 1:length(original)
            annotate!(p, orig_x[i] + offset, orig_y[i] + offset, orig_z[i] + offset, 
                    text("$(labels[i])", :red, 10, :right, 
                        strokecolor=:white, strokewidth=1, backgroundcolor=RGBA(1,1,1,0.7)))
            annotate!(p, trans_x[i] - offset, trans_y[i] - offset, trans_z[i] - offset, 
                    text("$(labels[i])'", :blue, 10, :left,
                        strokecolor=:white, strokewidth=1, backgroundcolor=RGBA(1,1,1,0.7)))
        end

        # Ось винта
        plot!(p, [0, 0], [0, 0], [0, 2], color=:green, linewidth=3, arrow=true, label="Ось (Oz)")
        
        return p
    end

    begin

    angle =  2π
    translation_distance = 1.0

    angle_slider = @bind angle Slider(0:0.01:2π, default=π, show_value=true)
    trans_slider = @bind translation_distance Slider(0:0.01:1.0, default=1.0, show_value=true)

    md"""
    ### Управление винтовым движением
    Угол поворота (рад): $(angle_slider)  
    Трансляция вдоль оси: $(trans_slider)
    """
    End

    begin
    # Ось винта — фиксирована (Oz)
    axis_direction = normalize([0.0, 0.0, 1.0])
    A = DualQuaternion(
        Quaternion(0.0, axis_direction...), 
        Quaternion(0.0, 0.0, 0.0, 0.0)
    )
    
    θ = Dual(angle, translation_distance)
    motor_func = PointSandwichFormula(A, θ)
    
    original_points = copy(points)
    transformed_points = []
    
    for p in points
        point_dq = point_to_dualquaternion(p)
        transformed_dq = motor_func(point_dq)
        transformed_p = dualquaternion_to_point(transformed_dq)
        push!(transformed_points, transformed_p)
    end
    
    plot_pyramids(original_points, transformed_points)
    end

\end{minted}

\section{Приложение №2}

\begin{minted}{Julia}
    begin
    using Quaternions
    using DualNumbers
    using LinearAlgebra
    
    # Вспомогательные функции для кватернионов
    function dot(q1::Quaternion, q2::Quaternion)
        return q1.s * q2.s + q1.v1 * q2.v1 + q1.v2 * q2.v2 + q1.v3 * q2.v3
    end
    
    function cross(q1::Quaternion, q2::Quaternion)
        # Векторное произведение для чисто векторных кватернионов
        return Quaternion(0.0, 
            q1.v2 * q2.v3 - q1.v3 * q2.v2,
            q1.v3 * q2.v1 - q1.v1 * q2.v3,
            q1.v1 * q2.v2 - q1.v2 * q2.v1
        )
    end
    
    function conj(q::Quaternion)
        return Quaternion(q.s, -q.v1, -q.v2, -q.v3)
    end
    
    function norm(q::Quaternion)
        return sqrt(dot(q, q))
    end
    
    struct DualQuaternion
        q::Quaternion{Float64}   # вещественная часть
        qo::Quaternion{Float64}  # дуальная часть
        
        # Конструктор из 4 дуальных чисел
        function DualQuaternion(Q0::Dual, Q1::Dual, Q2::Dual, Q3::Dual)
            q = Quaternion(Q0.value, Q1.value, Q2.value, Q3.value)
            qo = Quaternion(Q0.epsilon, Q1.epsilon, Q2.epsilon, Q3.epsilon)
            new(q, qo)
        end
        
        # Конструктор из двух кватернионов
        DualQuaternion(q::Quaternion, qo::Quaternion) = new(q, qo)
    end
    
    # Арифметические операции
    Base.:+(Q::DualQuaternion, P::DualQuaternion) = DualQuaternion(Q.q + P.q, Q.qo + P.qo)
    Base.:-(Q::DualQuaternion) = DualQuaternion(-Q.q, -Q.qo)
    Base.:-(Q::DualQuaternion, P::DualQuaternion) = DualQuaternion(Q.q - P.q, Q.qo - P.qo)
    
    # Сложение с дуальным числом (только скалярная часть)
    Base.:+(Q::DualQuaternion, d::Dual) = DualQuaternion(
        Quaternion(Q.q.s + d.value, Q.q.v1, Q.q.v2, Q.q.v3),
        Quaternion(Q.qo.s + d.epsilon, Q.qo.v1, Q.qo.v2, Q.qo.v3)
    )
    Base.:+(d::Dual, Q::DualQuaternion) = Q + d
    Base.:-(Q::DualQuaternion, d::Dual) = Q + (-d)
    Base.:-(d::Dual, Q::DualQuaternion) = DualQuaternion(
        Quaternion(d.value - Q.q.s, -Q.q.v1, -Q.q.v2, -Q.q.v3),
        Quaternion(d.epsilon - Q.qo.s, -Q.qo.v1, -Q.qo.v2, -Q.qo.v3)
    )
    
    # Умножение на число
    Base.:*(Q::DualQuaternion, a::Real) = DualQuaternion(a * Q.q, a * Q.qo)
    Base.:*(a::Real, Q::DualQuaternion) = Q * a
    
    # Умножение на дуальное число
    Base.:*(Q::DualQuaternion, d::Dual) = DualQuaternion(
        d.value * Q.q,
        d.value * Q.qo + d.epsilon * Q.q
    )
    Base.:*(d::Dual, Q::DualQuaternion) = Q * d
    
    # Умножение двух дуальных кватернионов
    Base.:*(Q::DualQuaternion, P::DualQuaternion) = DualQuaternion(
        Q.q * P.q,
        Q.q * P.qo + Q.qo * P.q
    )
    
    # Деление
    Base.:/(Q::DualQuaternion, a::Real) = DualQuaternion(Q.q / a, Q.qo / a)
    Base.:/(Q::DualQuaternion, d::Dual) = Q * inv(d)
    
    # Сравнение
    Base.:(==)(Q1::DualQuaternion, Q2::DualQuaternion) = (Q1.q == Q2.q) && (Q1.qo == Q2.qo)
    
    function Base.isapprox(Q1::DualQuaternion, Q2::DualQuaternion; 
                        atol::Real=1e-12, rtol::Real=0)
        return isapprox(Q1.q.s, Q2.q.s; atol, rtol) &&
                isapprox(Q1.q.v1, Q2.q.v1; atol, rtol) &&
                isapprox(Q1.q.v2, Q2.q.v2; atol, rtol) &&
                isapprox(Q1.q.v3, Q2.q.v3; atol, rtol) &&
                isapprox(Q1.qo.s, Q2.qo.s; atol, rtol) &&
                isapprox(Q1.qo.v1, Q2.qo.v1; atol, rtol) &&
                isapprox(Q1.qo.v2, Q2.qo.v2; atol, rtol) &&
                isapprox(Q1.qo.v3, Q2.qo.v3; atol, rtol)
    end
    
    # Сопряжения
    dualconj(Q::DualQuaternion) = DualQuaternion(Q.q, -Q.qo)
    complexconj(Q::DualQuaternion) = DualQuaternion(conj(Q.q), conj(Q.qo))
    conj(Q::DualQuaternion) = DualQuaternion(conj(Q.q), -conj(Q.qo))
    
    # Нормы и инварианты
    abs2(Q::DualQuaternion) = Dual(dot(Q.q, Q.q), 2 * dot(Q.q, Q.qo))
    invariant(Q::DualQuaternion) = dot(Q.q, Q.qo)
    parameter(Q::DualQuaternion) = invariant(Q) / dot(Q.q, Q.q)
    abs(Q::DualQuaternion) = Dual(norm(Q.q), invariant(Q) / norm(Q.q))
    
    # Конструкторы специальных объектов
    function screw(v::AbstractVector, m::AbstractVector)
        @id141972837 (@assert) length(v) == 3 && length(m) == 3 "Vectors must have length 3"
        return DualQuaternion(
            Quaternion(0.0, v[1], v[2], v[3]),
            Quaternion(0.0, m[1], m[2], m[3])
        )
    end
    
    function mass_point(w::Real, p::AbstractVector)
        @id141972837 (@assert) length(p) ==
    3 "Point must have length 3"
        return DualQuaternion(
            Quaternion(w, 0.0, 0.0, 0.0),
            Quaternion(0.0, p[1], p[2], p[3])
        )
    end
    
    function dq_plane(n::AbstractVector, d::Real)
        @id141972837 (@assert) length(n) == 3 "Normal vector must have length 3"
        return DualQuaternion(
            Quaternion(0.0, n[1], n[2], n[3]),
            Quaternion(d, 0.0, 0.0, 0.0)
        )
    end
    
    # Скалярное и векторное произведение
    dot(P::DualQuaternion, Q::DualQuaternion) = Dual(
        dot(P.q, Q.q),
        dot(P.qo, Q.q) + dot(P.q, Q.qo)
    )
    
    function cross(P::DualQuaternion, Q::DualQuaternion)
        @id141972837 (@assert) isapprox(P.q.s, 0.0) && isapprox(P.qo.s, 0.0) && 
                isapprox(Q.q.s, 0.0) && isapprox(Q.qo.s, 0.0) "Screw product requires pure dual quaternions"
        return DualQuaternion(
            cross(P.q, Q.q),
            cross(P.q, Q.qo) + cross(P.qo, Q.q)
        )
    end
    
    # Формулы для винтового движения
    function RodriguesHamilton(A::DualQuaternion, Θ::Dual)
        Λ0 = cos(Θ/2)
        Λ = sin(Θ/2) * A
        return L -> L + 2Λ0 * cross(Λ, L) + 2 * cross(Λ, cross(Λ, L))
    end
    
    function LineSandwichFormula(A::DualQuaternion, Θ::Dual)
        Λ = cos(Θ/2) + sin(Θ/2) * A
        return L -> Λ * L * complexconj(Λ)
    end
    
    function PointSandwichFormula(A::DualQuaternion, Θ::Dual)
        Λ = cos(Θ/2) + sin(Θ/2) * A
        return P -> Λ * P * conj(Λ)
    end
    
    PlaneSandwichFormula = PointSandwichFormula
    
    # Вспомогательная функция для ортогональных векторов
    function orthogonal_triple(n::AbstractVector)
        @id141972837 (@assert) length(n) == 3 "Vector must have length 3"
        if n[1] <= n[2] && n[1] <= n[3]
            w = n + [1.0, 0.0, 0.0]
        elseif n[2] <= n[1] && n[2] <= n[3]
            w = n + [0.0, 1.0, 0.0]
        else
            w = n + [0.0, 0.0, 1.0]
        end
        u = normalize(cross(w, n))
        v = normalize(cross(n, u))
        return (v, u)
    end
    
    
    end

    begin
    using Plots
    using Colors
        using PlutoUI
    end

    begin

        angle = 2π
        translation_distance = 1.0

        angle_slider = @bind angle Slider(0:0.01:2π, default=π, show_value=true)
        trans_slider = @bind translation_distance Slider(0:0.01:1.0, default=1.0, show_value=true)

        md"""
        ### Управление винтовым движением
        Угол поворота (рад): $(angle_slider)  
        Трансляция вдоль оси: $(trans_slider)
        """
    end

    begin
        points = [
            [0.0, 0.0, 0.0],      # A
            [1.0, 0.0, 0.0],      # B
            [1.0, 1.0, 0.0],      # C
            [0.0, 1.0, 0.0],      # D
            [0.5, 0.5, √2 / 2]    # S
        ]

        base_faces = [[1,2,3,4]]
        side_faces = [[1,2,5], [2,3,5], [3,4,5], [4,1,5]]

    ъ    function point_to_dualquaternion(p::Vector)
            return DualQuaternion(
                Quaternion(1.0, 0.0, 0.0, 0.0),
                Quaternion(0.0, p[1], p[2], p[3])
            )
        end

        function dualquaternion_to_point(dq::DualQuaternion)
            return [dq.qo.v1, dq.qo.v2, dq.qo.v3]
        end

        function plot_pyramids_with_arbitrary_axis(original, transformed, axis_vector)
            p = plot(size=(1000, 800), legend=:topleft, 
                    xlabel="X", ylabel="Y", zlabel="Z",
                    camera=(30, 30),
                    xlims=(-2, 2), ylims=(-2, 2), zlims=(-0.2, 2.2))

            orig_x = [pt[1] for pt in original]
            orig_y = [pt[2] for pt in original] 
            orig_z = [pt[3] for pt in original]

            trans_x = [pt[1] for pt in transformed]
            trans_y = [pt[2] for pt in transformed]
            trans_z = [pt[3] for pt in transformed]

            base_x = [orig_x[1], orig_x[2], orig_x[3], orig_x[4], orig_x[1]]
            base_y = [orig_y[1], orig_y[2], orig_y[3], orig_y[4], orig_y[1]]
            base_z = [orig_z[1], orig_z[2], orig_z[3], orig_z[4], orig_z[1]]
            plot!(p, base_x, base_y, base_z, color=:red, linewidth=2, label="Исходная")

            for face in side_faces
                fx = [orig_x[i] for i in face]
                fy = [orig_y[i] for i in face] 
                fz = [orig_z[i] for i in face]
                plot!(p, fx,
    fy, fz, color=:red, linewidth=2, alpha=0.7, label="")
            end

            tbase_x = [trans_x[1], trans_x[2], trans_x[3], trans_x[4], trans_x[1]]
            tbase_y = [trans_y[1], trans_y[2], trans_y[3], trans_y[4], trans_y[1]]
            tbase_z = [trans_z[1], trans_z[2], trans_z[3], trans_z[4], trans_z[1]]
            plot!(p, tbase_x, tbase_y, tbase_z, color=:blue, linewidth=2, label="После движения")

            for face in side_faces
                fx = [trans_x[i] for i in face]
                fy = [trans_y[i] for i in face]
                fz = [trans_z[i] for i in face]
                plot!(p, fx, fy, fz, color=:blue, linewidth=2, alpha=0.7, label="")
            end

            scatter!(p, orig_x, orig_y, orig_z, color=:red, markersize=6, label="")
            scatter!(p, trans_x, trans_y, trans_z, color=:blue, markersize=6, label="")

            labels = ["A", "B", "C", "D", "S"]
            offset = 0.05
            for i in 1:length(original)
                annotate!(p, orig_x[i] + offset, orig_y[i] + offset, orig_z[i] + offset, 
                        text("$(labels[i])", :red, 10, :right, 
                            strokecolor=:white, strokewidth=2, backgroundcolor=RGBA(1,1,1,0.7)))
                annotate!(p, trans_x[i] - offset, trans_y[i] - offset, trans_z[i] - offset, 
                        text("$(labels[i])'", :blue, 10, :left,
                            strokecolor=:white, strokewidth=2, backgroundcolor=RGBA(1,1,1,0.7)))
            end

            # Ось винта
            axis_end = axis_vector .* 2.0
            plot!(p, [0, axis_end[1]], [0, axis_end[2]], [0, axis_end[3]], 
                color=:green, linewidth=3, arrow=true, label="Ось винта")

            return p
        end

        function spherical_to_cartesian(θ_deg::Float64, φ_deg::Float64)
            θ = deg2rad(θ_deg)
            φ = deg2rad(φ_deg)
            normalize([sin(θ)[club172098281|φ], sin(θ)φ, cos(θ)])
        end

        polar_angle = 60.0    # от оси Z
        azimuth_angle = 45.0  # от оси X
        axis_direction = spherical_to_cartesian(polar_angle, azimuth_angle)

        A = DualQuaternion(
            Quaternion(0.0, axis_direction...), 
            Quaternion(0.0, 0.0, 0.0, 0.0)
        )
        θ_dual = Dual(angle, translation_distance)  
        motor_func = PointSandwichFormula(A, θ_dual)

        # === Преобразование точек ===
        original_points = copy(points)
        transformed_points = [
            dualquaternion_to_point(motor_func(point_to_dualquaternion(p)))
            for p in points
        ]
    savefig("2.pdf")
        plot_pyramids_with_arbitrary_axis(original_points, transformed_points, axis_direction)

    
    end
\end{minted}

\section{Приложение №3}

\begin{minted}{Julia}
    begin
        using Quaternions
        using DualNumbers
        using LinearAlgebra
        
        # Вспомогательные функции для кватернионов
        function dot(q1::Quaternion, q2::Quaternion)
            return q1.s * q2.s + q1.v1 * q2.v1 + q1.v2 * q2.v2 + q1.v3 * q2.v3
        end
        
        function cross(q1::Quaternion, q2::Quaternion)
            # Векторное произведение для чисто векторных кватернионов
            return Quaternion(0.0, 
                q1.v2 * q2.v3 - q1.v3 * q2.v2,
                q1.v3 * q2.v1 - q1.v1 * q2.v3,
                q1.v1 * q2.v2 - q1.v2 * q2.v1
            )
        end
        
        function conj(q::Quaternion)
            return Quaternion(q.s, -q.v1, -q.v2, -q.v3)
        end
        
        function norm(q::Quaternion)
            return sqrt(dot(q, q))
        end
        
        struct DualQuaternion
            q::Quaternion{Float64}   # вещественная часть
            qo::Quaternion{Float64}  # дуальная часть
            
            # Конструктор из 4 дуальных чисел
            function DualQuaternion(Q0::Dual, Q1::Dual, Q2::Dual, Q3::Dual)
                q = Quaternion(Q0.value, Q1.value, Q2.value, Q3.value)
                qo = Quaternion(Q0.epsilon, Q1.epsilon, Q2.epsilon, Q3.epsilon)
                new(q, qo)
            end
            
            # Конструктор из двух кватернионов
            DualQuaternion(q::Quaternion, qo::Quaternion) = new(q, qo)
        end
        
        # Арифметические операции
        Base.:+(Q::DualQuaternion, P::DualQuaternion) = DualQuaternion(Q.q + P.q, Q.qo + P.qo)
        Base.:-(Q::DualQuaternion) = DualQuaternion(-Q.q, -Q.qo)
        Base.:-(Q::DualQuaternion, P::DualQuaternion) = DualQuaternion(Q.q - P.q, Q.qo - P.qo)
        
        # Сложение с дуальным числом (только скалярная часть)
        Base.:+(Q::DualQuaternion, d::Dual) = DualQuaternion(
            Quaternion(Q.q.s + d.value, Q.q.v1, Q.q.v2, Q.q.v3),
            Quaternion(Q.qo.s + d.epsilon, Q.qo.v1, Q.qo.v2, Q.qo.v3)
        )
        Base.:+(d::Dual, Q::DualQuaternion) = Q + d
        Base.:-(Q::DualQuaternion, d::Dual) = Q + (-d)
        Base.:-(d::Dual, Q::DualQuaternion) = DualQuaternion(
            Quaternion(d.value - Q.q.s, -Q.q.v1, -Q.q.v2, -Q.q.v3),
            Quaternion(d.epsilon - Q.qo.s, -Q.qo.v1, -Q.qo.v2, -Q.qo.v3)
        )
        
        # Умножение на число
        Base.:*(Q::DualQuaternion, a::Real) = DualQuaternion(a * Q.q, a * Q.qo)
        Base.:*(a::Real, Q::DualQuaternion) = Q * a
        
        # Умножение на дуальное число
        Base.:*(Q::DualQuaternion, d::Dual) = DualQuaternion(
            d.value * Q.q,
            d.value * Q.qo + d.epsilon * Q.q
        )
        Base.:*(d::Dual, Q::DualQuaternion) = Q * d
        
        # Умножение двух дуальных кватернионов
        Base.:*(Q::DualQuaternion, P::DualQuaternion) = DualQuaternion(
            Q.q * P.q,
            Q.q * P.qo + Q.qo * P.q
        )
        
        # Деление
        Base.:/(Q::DualQuaternion, a::Real) = DualQuaternion(Q.q / a, Q.qo / a)
        Base.:/(Q::DualQuaternion, d::Dual) = Q * inv(d)
        
        # Сравнение
        Base.:(==)(Q1::DualQuaternion, Q2::DualQuaternion) = (Q1.q == Q2.q) && (Q1.qo == Q2.qo)
        
        function Base.isapprox(Q1::DualQuaternion, Q2::DualQuaternion; 
                              atol::Real=1e-12, rtol::Real=0)
            return isapprox(Q1.q.s, Q2.q.s; atol, rtol) &&
                   isapprox(Q1.q.v1, Q2.q.v1; atol, rtol) &&
                   isapprox(Q1.q.v2, Q2.q.v2; atol, rtol) &&
                   isapprox(Q1.q.v3, Q2.q.v3; atol, rtol) &&
                   isapprox(Q1.qo.s, Q2.qo.s; atol, rtol) &&
                   isapprox(Q1.qo.v1, Q2.qo.v1; atol, rtol) &&
                   isapprox(Q1.qo.v2, Q2.qo.v2; atol, rtol) &&
                   isapprox(Q1.qo.v3, Q2.qo.v3; atol, rtol)
        end
        
        # Сопряжения
        dualconj(Q::DualQuaternion) = DualQuaternion(Q.q, -Q.qo)
        complexconj(Q::DualQuaternion) = DualQuaternion(conj(Q.q), conj(Q.qo))
        conj(Q::DualQuaternion) = DualQuaternion(conj(Q.q), -conj(Q.qo))
        
        # Нормы и инварианты
        abs2(Q::DualQuaternion) = Dual(dot(Q.q, Q.q), 2 * dot(Q.q, Q.qo))
        invariant(Q::DualQuaternion) = dot(Q.q, Q.qo)
        parameter(Q::DualQuaternion) = invariant(Q) / dot(Q.q, Q.q)
        abs(Q::DualQuaternion) = Dual(norm(Q.q), invariant(Q) / norm(Q.q))
        
        # Конструкторы специальных объектов
        function screw(v::AbstractVector, m::AbstractVector)
            @assert length(v) == 3 && length(m) == 3 "Vectors must have length 3"
            return DualQuaternion(
                Quaternion(0.0, v[1], v[2], v[3]),
                Quaternion(0.0, m[1], m[2], m[3])
            )
        end
        
        function mass_point(w::Real, p::AbstractVector)
            @assert length(p) == 3 "Point must have length 3"
            return DualQuaternion(
                Quaternion(w, 0.0, 0.0, 0.0),
                Quaternion(0.0, p[1], p[2], p[3])
            )
        end
        
        function dq_plane(n::AbstractVector, d::Real)
            @assert length(n) == 3 "Normal vector must have length 3"
            return DualQuaternion(
                Quaternion(0.0, n[1], n[2], n[3]),
                Quaternion(d, 0.0, 0.0, 0.0)
            )
        end
        
        # Скалярное и векторное произведение
        dot(P::DualQuaternion, Q::DualQuaternion) = Dual(
            dot(P.q, Q.q),
            dot(P.qo, Q.q) + dot(P.q, Q.qo)
        )
        
        function cross(P::DualQuaternion, Q::DualQuaternion)
            @assert isapprox(P.q.s, 0.0) && isapprox(P.qo.s, 0.0) && 
                    isapprox(Q.q.s, 0.0) && isapprox(Q.qo.s, 0.0) "Screw product requires pure dual quaternions"
            return DualQuaternion(
                cross(P.q, Q.q),
                cross(P.q, Q.qo) + cross(P.qo, Q.q)
            )
        end
        
        # Формулы для винтового движения
        function RodriguesHamilton(A::DualQuaternion, Θ::Dual)
            Λ0 = cos(Θ/2)
            Λ = sin(Θ/2) * A
            return L -> L + 2Λ0 * cross(Λ, L) + 2 * cross(Λ, cross(Λ, L))
        end
        
        function LineSandwichFormula(A::DualQuaternion, Θ::Dual)
            Λ = cos(Θ/2) + sin(Θ/2) * A
            return L -> Λ * L * complexconj(Λ)
        end
        
        function PointSandwichFormula(A::DualQuaternion, Θ::Dual)
            Λ = cos(Θ/2) + sin(Θ/2) * A
            return P -> Λ * P * conj(Λ)
        end
        
        PlaneSandwichFormula = PointSandwichFormula
        
        # Вспомогательная функция для ортогональных векторов
        function orthogonal_triple(n::AbstractVector)
            @assert length(n) == 3 "Vector must have length 3"
            if n[1] <= n[2] && n[1] <= n[3]
                w = n + [1.0, 0.0, 0.0]
            elseif n[2] <= n[1] && n[2] <= n[3]
                w = n + [0.0, 1.0, 0.0]
            else
                w = n + [0.0, 0.0, 1.0]
            end
            u = normalize(cross(w, n))
            v = normalize(cross(n, u))
            return (v, u)
        end
        
        
    end
    begin
    # Определение точки P (аффинная точка)
    point_coords = [0.5, 0.5, √2 / 2]
    m = DualQuaternion(Quaternion(1.0, 0.0, 0.0, 0.0), Quaternion(0.0, point_coords...)) # m=1+(xi+yj+zk)ε
    end

    begin
    # Разделы 5.4 и 5.3
    # Определение оси вращения (проходящую через начало координат)
    # Например, ось Z: s = (0, 0, 1)
    axis_direction = normalize([0.0, 0.0, 1.0]) # Проверить, что нормирован
    # Момент для оси через O: s_0 = (0, 0, 0)
    A = DualQuaternion(Quaternion(0.0, axis_direction...), Quaternion(0.0, 0.0, 0.0, 0.0))
    # Должно удовлетворяться условие Плюккера
    end

    begin
    # На 90 градусов
    angle = π/2
    Θ = Dual(angle, 0.0)
    end

    #  Бикватернион винтового движения определен в блоке выше(формула 12, стр 19)
    # Предполагается использование сэндвич-формулы(раздел 5.4.1), бикватернионное сопряжение раздел 5.2.5
    rot_func = PointSandwichFormula(A, Θ)

    begin
    # Применяем вращение
    m_rotated = rot_func(m)
    
    println("Исходное представление точки: ", m)
    println("Представление после вращения: ", m_rotated)
    end

    begin
    # Извлекаем координаты вращеной точки из m_rotated.qo (векторная часть дуальной части)
    # потому что m_rotated должно быть вида (1, 0, 0, 0) + (0, x', y', z')*ε
    # для аффинной точки (x', y', z', 1)
    rotated_point_coords = [m_rotated.qo.v1, m_rotated.qo.v2, m_rotated.qo.v3]
    println("Исходные координаты точки: ", point_coords)
    println("Координаты точки после вращения: ", rotated_point_coords)
    
    end

    begin
    using Plots
    using Colors
    end

    points = [
        [0.0, 0.0, 0.0],      # A
        [1.0, 0.0, 0.0],      # B
        [1.0, 1.0, 0.0],      # C
        [0.0, 1.0, 0.0],      # D
        [0.5, 0.5, √2 / 2]    # S (вершина)
    ]

    # Функция для преобразования точки в бикватернион
    function point_to_dualquaternion(p::Vector)
        return DualQuaternion(
            Quaternion(1.0, 0.0, 0.0, 0.0),  
            Quaternion(0.0, p[1], p[2], p[3]) 
        )
    end

    # Функция для извлечения координат из бикватерниона точки
    function dualquaternion_to_point(dq::DualQuaternion)
        return [dq.qo.v1, dq.qo.v2, dq.qo.v3]
    end

    begin
    # Параметры винтового движения
    angle1 = π  
    translation_distance = 0.0  
    θ1 = Dual(angle1, translation_distance)  
    end

    begin
    # Ось винтового движения (Oz)
    axis_direction1 = normalize([0.0, 0.0, 1.0])
    A1 = DualQuaternion(
        Quaternion(0.0, axis_direction1...), 
        Quaternion(0.0, 0.0, 0.0, 0.0)
    )
    end

    motor_func = PointSandwichFormula(A1, θ1)

    begin
    # Применяем винтовое движение ко всем точкам
    original_points = []
    transformed_points = []
    
    for p in points
        # Преобразуем точку в бикватернион
        point_dq = point_to_dualquaternion(p)
        
        # Применяем винтовое движение
        transformed_dq = motor_func(point_dq)
        
        # Извлекаем координаты
        transformed_p = dualquaternion_to_point(transformed_dq)
        
        push!(original_points, p)
        push!(transformed_points, transformed_p)
    end
    end

    begin
    # Основание ABCD
    base_faces = [[1,2,3,4]]  # A-B-C-D
    
    # Боковые грани
    side_faces = [
        [1,2,5],  # A-B-S
        [2,3,5],  # B-C-S  
        [3,4,5],  # C-D-S
        [4,1,5]   # D-A-S
    ]
    
    end

    begin
        normal_vec = [nx, ny, nz]
        n_len = LinearAlgebra.norm(normal_vec)
        
        if n_len > 0 
            normal_vec = normal_vec / n_len 
        else 
            normal_vec = [0.0, 0.0, 1.0] 
        end

        reflection_plane = dq_plane(normal_vec, d_plane)

    reflected_points = []
    for p in transformed_points
        dq = point_to_dualquaternion(p)
        # Используйте правильную формулу
        dq_reflected = reflect_point(dq, reflection_plane)
        p_reflected = dualquaternion_to_point(dq_reflected)
        push!(reflected_points, p_reflected)
    end
    
    end

    using Statistics

    function plot_pyramids(original, transformed, reflected, normal, d)

        
        p = plot(size=(1000, 800), legend=:topleft, xlabel="x", ylabel="y", zlabel="Z", 
                camera=(30, 30), xlims=(-2.5, 2.5), ylims=(-2.5, 2.5), zlims=(-2.5, 2.5),
                grid=true, aspect_ratio=:equal)

        center_x = mean([pt[1] for pt in transformed])
        center_y = mean([pt[2] for pt in transformed])
        center_z = mean([pt[3] for pt in transformed])
        center = [center_x, center_y,
    center_z]

        n_dot_c = LinearAlgebra.dot(normal, center)
        dist = n_dot_c - d
        center_on_plane = center - dist .* normal

        range_size = 2.5
        
        xs = range(center_on_plane[1] - range_size, center_on_plane[1] + range_size, length=30)
        ys = range(center_on_plane[2] - range_size, center_on_plane[2] + range_size, length=30)

        if abs(normal[3]) > 1e-6
            zs = [(d - normal[1]*x - normal[2]*y) / normal[3] for x in xs, y in ys]
            zs = clamp.(zs, -2.5, 2.5)
            surface!(p, xs, ys, zs, color=:gray, alpha=0.3, legend=false)
        elseif abs(normal[2]) > 1e-6
            # Если плоскость вертикальна вдоль Y
            zs = range(center_on_plane[3] - range_size, center_on_plane[3] + range_size, length=30)
            ys_calc = [(d - normal[1]*x - normal[3]*z) / normal[2] for x in xs, z in zs]
            ys_calc = clamp.(ys_calc, -2.5, 2.5)
            surface!(p, xs, ys_calc, zs, color=:gray, alpha=0.3, legend=false)
        else
            # Если плоскость вертикальна вдоль X
            zs = range(center_on_plane[3] - range_size, center_on_plane[3] + range_size, length=30)
            ys_calc = range(center_on_plane[2] - range_size, center_on_plane[2] + range_size, length=30)
            xs_calc = [(d - normal[2]*y - normal[3]*z) / normal[1] for y in ys_calc, z in zs]
            xs_calc = clamp.(xs_calc, -2.5, 2.5)
            surface!(p, xs_calc, ys_calc, zs, color=:gray, alpha=0.3, legend=false)
        end

        plot!(p, [0, 0], [0, 0], [-2.5, 2.5], color=:green, linewidth=3, arrow=true, label="Ось винта")

        orig_x = [pt[1] for pt in original]
        orig_y = [pt[2] for pt in original]
        orig_z = [pt[3] for pt in original]
        
        base_x = [orig_x[1], orig_x[2], orig_x[3], orig_x[4], orig_x[1]]
        base_y = [orig_y[1], orig_y[2], orig_y[3], orig_y[4], orig_y[1]]
        base_z = [orig_z[1], orig_z[2], orig_z[3], orig_z[4], orig_z[1]]
        
        plot!(p, base_x, base_y, base_z, color=:red, linewidth=2, label="Исходная")
        
        side_faces = [[1,2,5],[2,3,5],[3,4,5],[4,1,5]]
        for face in side_faces
            fx = [orig_x[i] for i in face]
            fy = [orig_y[i] for i in face]
            fz = [orig_z[i] for i in face]
            plot!(p, fx, fy, fz, color=:red, linewidth=2, alpha=0.7, label="")
        end
        scatter!(p, orig_x, orig_y, orig_z, color=:red, markersize=6, label="")

        trans_x = [pt[1] for pt in transformed]
        trans_y = [pt[2] for pt in transformed]
        trans_z = [pt[3] for pt in transformed]
        
        trans_base_x = [trans_x[1], trans_x[2], trans_x[3], trans_x[4], trans_x[1]]
        trans_base_y = [trans_y[1], trans_y[2], trans_y[3], trans_y[4], trans_y[1]]
        trans_base_z = [trans_z[1], trans_z[2], trans_z[3], trans_z[4], trans_z[1]]
        
        plot!(p, trans_base_x, trans_base_y, trans_base_z, color=:blue, linewidth=2, label="После винта")
        
        for face in side_faces
            fx = [trans_x[i] for i in face]
            fy = [trans_y[i] for i in face]
            fz = [trans_z[i] for i in face]
            plot!(p, fx, fy, fz, color=:blue, linewidth=2, alpha=0.7, label="")
        end
        scatter!(p, trans_x, trans_y, trans_z, color=:blue, markersize=6, label="")

        refl_x = [pt[1] for pt in reflected]
        refl_y = [pt[2] for pt in reflected]
        refl_z = [pt[3] for pt in reflected]
        
        refl_base_x = [refl_x[1], refl_x[2], refl_x[3], refl_x[4], refl_x[1]]
        refl_base_y = [refl_y[1], refl_y[2], refl_y[3], refl_y[4], refl_y[1]]
        refl_base_z = [refl_z[1], refl_z[2], refl_z[3], refl_z[4], refl_z[1]]
        
        plot!(p, refl_base_x, refl_base_y, refl_base_z, color=:violet, linewidth=2, label="Отражённая")
        
        for face in side_faces
            fx = [refl_x[i] for i in face]
            fy = [refl_y[i] for i in face]
            fz = [refl_z[i] for i in face]
            plot!(p, fx, fy, fz, color=:violet, linewidth=2, alpha=0.7, label="")
        end
        scatter!(p, refl_x, refl_y, refl_z, color=:violet, markersize=6, label="")

        labels = ["A","B","C","D","S"]
        offset = 0.12
        for i in 1:length(original)
            annotate!(p, orig_x[i]+offset, orig_y[i]+offset, orig_z[i]+offset, text("$(labels[i])", :red, 12, :right, strokecolor=:white, strokewidth=2))
    annotate!(p, trans_x[i]-offset, trans_y[i]-offset, trans_z[i]-offset, text("$(labels[i])", :blue, 12, :left, strokecolor=:white, strokewidth=2))
            annotate!(p, refl_x[i], refl_y[i]+offset, refl_z[i]+offset, text("$(labels[i])", :violet, 12, :center, strokecolor=:white, strokewidth=2))
        end
        
        return p
    end

    begin
        @bind nx Slider(-1:0.01:1, default=0.0, show_value=true)
        @bind ny Slider(-1:0.01:1, default=0.0, show_value=true)
        @bind nz Slider(-1:0.01:1, default=1.0, show_value=true)
        @bind d_plane Slider(-2:0.1:2, default=0.5, show_value=true)
        
        md"""
        ### Настройка плоскости отражения
        
        **Нормальный вектор:**
        - nx: $(@bind nx Slider(-1:0.01:1, default=0.0, show_value=true))
        - ny: $(@bind ny Slider(-1:0.01:1, default=0.0, show_value=true))
        - nz: $(@bind nz Slider(-1:0.01:1, default=1.0, show_value=true))
        
        **Расстояние от начала координат:**
        - d: $(@bind d_plane Slider(-2:0.1:2, default=0.5, show_value=true))
        
        """
    end

    begin
        plot_result = plot_pyramids(original_points, transformed_points, reflected_points, normal_vec, d_plane)

    savefig(plot_result, "пирамида_преобразование.pdf")
        
        println("Исходные точки пирамиды:")
        labels = ["A", "B", "C", "D", "S"]
        for (i, p) in enumerate(original_points)
            println("$(labels[i]): $p")
        end
        
        println("\nПосле винтового движения (вращение 180° вокруг Oz + трансляция на 1 вдоль Oz):")
        for (i, p) in enumerate(transformed_points)
            println("$(labels[i])': $p")
        end
        
        println("\nТекущая плоскость отражения: $(round.(normal_vec, digits=2)) · r = $(round(d_plane, digits=2))")
        
        plot_result
    end
    
\end{minted}

\section{Приложение №4}

\begin{minted}{python}
    import matplotlib.pyplot as plt
    import numpy as np
    import math

    from math import cos, sin

    from dataclasses import dataclass



    @dataclass
    class Dual:
        real : float       
        dual : float = 0   
            
        def __add__(self, y):
            if not isinstance(y, Dual):
                y = Dual(y, 0)
            return Dual(self.real + y.real, self.dual + y.dual)
        
        def __sub__(self, y):
            if not isinstance(y, Dual):
                y = Dual(y, 0)        
            return Dual(self.real - y.real, self.dual - y.dual)
        
        def __mul__(self, y):
            if not isinstance(y, Dual):
                y = Dual(y, 0)        
            return Dual(self.real * y.real, (self.real * y.dual) + (self.dual * y.real))
        
        def __truediv__(self, y):
            if not isinstance(y, Dual):
                y = Dual(y, 0)        
            y_real_inv = 1 / y.real
            real_div = self.real * y_real_inv
            return Dual(real_div, (self.dual - real_div*y.dual) * y_real_inv)

        def __repr__(self):
            if self.dual >= 0:
                return f'{self.real} + {self.dual}ε'
            else: 
                return f'{self.real} + {-self.dual}ε'

        @staticmethod
        def cos(A):
            "Косинус от дуального числа"
            return Dual(cos(A.real), -sin(A.real) * A.dual)

        @staticmethod
        def sin(A):
            "Синус от дуального числа"
            return Dual(sin(A.real), cos(A.real) * A.dual)

    # Статический метод вызывается не из объекта, а из класса
    print(Dual.cos(Dual(math.pi/2, 1)))
    print(Dual.sin(Dual(math.pi/2, 1)))

    # Точка O (начало координат)  
    O = np.array([0, 0, 1])
    # Направляющий вектор прямой
    v = np.array([2, 1, 0])
    # Вычисление момента для прямой, проходящей через точку P  с направляющим вектором v
    p = O
    m = np.cross(p, v)
    # Определяем винт, путем задания его дуальных координат.
    # Действительную часть берем из вектора v, а дуальную часть из момента m
    L = np.array([Dual(v[i], m[i]) for i in range(3)], dtype=Dual)
    # Теперь L — винт, задающий прямую

    # Направляющий вектор оси Oz
    OZ = np.array([0, 0, 1])
    # Зададим теперь ось вращения в виде винта
    A = np.array([Dual(v, m) for (v, m) in zip(OZ, np.cross(O, OZ))], dtype=Dual)

    print("Вектор v:", v)
    print("Вектор p:", p)
    print("Вектор m (момент):", m)
    print(f"Винт L = v + ε * m = {L}")
    print(f"Винт A = v + ε * m = {A}")

    def Rodrig(A, Θ, L):
        return A * (Dual(1, 0) - Dual.cos(Θ)) * np.dot(A, L) + L * Dual.cos(Θ) + np.cross(A, L) * Dual.sin(Θ)

    result = Rodrig(A, Θ, L)
    print(f"Результат функции Rodrig:\n{result}")

    Θ = Dual(math.pi / 2, 1)  # дуальный угол

    # Вычисление нового вектора
    L1 = Rodrig(A, Θ, L)

    print(f"L1' = {L1}")

    real_parts = np.array([component.real for component in L1])
    dual_parts = np.array([component.dual for component in L1])

    print("Действительные части:", real_parts)
    print("Дуальные части:", dual_parts)


    #Соответственно  v' обозначим, как V
    V = real_parts
    V

    # m обозначим как  M
    M = dual_parts
    M

    # Вычисление нормы вектора v
    norm = np.linalg.norm(v)

    # Нормализация вектора
    normalized_v = v / norm

    print("Нормированный вектор:", normalized_v)


    # Соответствует [ 1/√2 1/√2 0] так,как 1√2 = 0.707...


    # Вычисление нормы вектора M
    norm = np.linalg.norm(M)

    # Нормализация вектора
    normalized_M = M / norm

    # Вывод результатов
    print("Нормированный вектор:", normalized_M)


    # Соответствует [ -1 0 0] так,как 7.85046229e-17 является очень маленьким числом, близким к нулю.


    # Вычисление нормы вектора V
    norm = np.linalg.norm(V)

    # Нормализация вектора
    normalized_V = V / norm

    # Вывод результатов
    print("Нормированный вектор:", normalized_V)

    # Соответствует [ 0 1 0] так,как 7.85046229e-17 является очень маленьким числом, близким к нулю.


    P = np.cross(V, M) / np.linalg.norm(V)**2
    P


    # Прямая L1 будет проходить через точку [0 0 -1] 


    fig = plt.figure(num=1, figsize=(8, 8), dpi=200)
    ax = fig.add_subplot(111, projection='3d')

    ax.set_xlim([-5, 5])
    ax.set_ylim([-5, 5])
    ax.set_zlim([-5, 5])

    axis_length = 5

    ax.plot([0, axis_length], [0, 0], [0, 0], color='red', linewidth=1.5, linestyle='-', zorder=3)  # X+
    ax.plot([0, 0], [0, axis_length], [0, 0], color='green', linewidth=1.5, linestyle='-', zorder=3)  # Y+
    ax.plot([0, 0], [0, 0], [0, axis_length], color='blue', linewidth=1.5, linestyle='-', zorder=3)  # Z+

    ax.plot([0, -axis_length], [0, 0], [0, 0], color='red', linewidth=1.5, linestyle='--', zorder=3)  # X-
    ax.plot([0, 0], [0, -axis_length], [0, 0], color='green', linewidth=1.5, linestyle='--', zorder=3)  # Y-
    ax.plot([0, 0], [0, 0], [0, -axis_length], color='blue', linewidth=1.5, linestyle='--', zorder=3)  # Z-

    ax.scatter(*O, s=30, c='k', label='Точка O (0, 0, 0)')
    ax.scatter(*(v + O), s=30, c='purple', label=f'Вектор прямой v')
    ax.scatter(*(P + O), s=30, c='red', label=f'Точка P')

    t = np.linspace(-5, 5, 100)  
    x = O[0] + v[0] * t
    y = O[1] + v[1] * t
    z = O[2] + v[2] * t

    ax.plot(x, y, z, label='L', linewidth=2)

    t2 = np.linspace(-5, 5, 100)
    x2 = P[0] + normalized_V[0] * t2 
    y2 = P[1] + normalized_V[1] * t2 
    z2 = P[2] + normalized_V[2] * t2 
    ax.plot(x2, y2, z2, label='L1', linewidth=2)


    ax.set_xlim([-5, 5])
    ax.set_ylim([-5, 5])
    ax.set_zlim([-5, 5])

    axis_length = 5

    ax.set_xlabel('X')
    ax.set_ylabel('Y')
    ax.set_zlabel('Z')

    ax.legend(fontsize=10)

    ax.view_init(elev=20, azim=30)

    plt.tight_layout()

    plt.savefig('graph.pdf', format='pdf')

    plt.show()

    from ipywidgets import FloatSlider, VBox, interactive
    from IPython.display import display

    def update_dual(real_part, dual_part):
        Θ = Dual(real_part, dual_part)  # дуальный угол
        L1 = Rodrig(A, Θ, L)  

        real_parts = np.array([component.real for component in L1])
        dual_parts = np.array([component.dual for component in L1])

        plt.clf()

        fig = plt.figure(num=1, figsize=(8, 8), dpi=200)
        ax = fig.add_subplot(111, projection='3d')

        ax.set_xlim([-5, 5])
        ax.set_ylim([-5, 5])
        ax.set_zlim([-5, 5])

        axis_length = 5

        ax.plot([0, axis_length], [0, 0], [0, 0], color='red', linewidth=1.5, linestyle='-', zorder=3)  # X+
        ax.plot([0, 0], [0, axis_length], [0, 0], color='green', linewidth=1.5, linestyle='-', zorder=3)  # Y+
        ax.plot([0, 0], [0, 0], [0, axis_length], color='blue', linewidth=1.5, linestyle='-', zorder=3)  # Z+

        ax.plot([0, -axis_length], [0, 0], [0, 0], color='red', linewidth=1.5, linestyle='--', zorder=3)  # X-
        ax.plot([0, 0], [0, -axis_length], [0, 0], color='green', linewidth=1.5, linestyle='--', zorder=3)  # Y-
        ax.plot([0, 0], [0, 0], [0, -axis_length], color='blue', linewidth=1.5, linestyle='--', zorder=3)  # Z-

        ax.scatter(*O, s=30, c='k', label='Точка O (0, 0, 0)')
        ax.scatter(*(v + O), s=30, c='purple', label='Вектор линии v')
        
        P = np.cross(real_parts, dual_parts) / np.linalg.norm(real_parts)**2
        ax.scatter(*(P + O), s=30, c='red', label='Точка P')

        t = np.linspace(-5, 5, 100)  
        x = O[0] + v[0] * t
        y = O[1] + v[1] * t
        z = O[2] + v[2] * t

        ax.plot(x, y, z, label='Прямая L', linewidth=2)

        t2 = np.linspace(-5, 5, 100)  
        x2 = P[0] + real_parts[0] * t2
        y2 = P[1] + real_parts[1] * t2
        z2 = P[2] + real_parts[2] * t2

        ax.plot(x2, y2, z2, label='Прямая L1', linewidth=2)

        ax.set_xlabel('X')
        ax.set_ylabel('Y')
        ax.set_zlabel('Z')
        ax.legend(fontsize=10)

        ax.view_init(elev=20, azim=30)
        plt.tight_layout()
        plt.show()

    real_slider = FloatSlider(value=Θ.real, min=0, max=math.pi / 4, step=0.01, description='Real part:')
    dual_slider = FloatSlider(value=Θ.dual, min=0, max=1, step=0.01, description='Dual part:')

    interactive_plot = interactive(update_dual, real_part=real_slider, dual_part=dual_slider)
    output = interactive_plot.children[-1]
    output.layout.height = '400px'

    display(real_slider, dual_slider, output)
\end{minted}

